% =====================================================================
%  Lecture 04 -- Mobility and Handover                    (Persian edition)
% =====================================================================
\documentclass[11pt,a4paper]{article}
\providecommand\ACNroot{../..}
\input{\ACNroot/templates/lecture-fa}

\ACNset{number}{04}
\ACNset{title}{تحرک و تحویل}
\ACNset{shorttitle}{تحرک و تحویل}
\ACNset{titleother}{Mobility and Handover}
\ACNset{part}{بخش اول -- مبانی نسل دوم و سوم (درس‌های ۱ تا ۴)}
\ACNset{time}{۷ تا ۹ ساعت (۳٫۵ ساعت مطالعه، ۴ تا ۵ ساعت تمرین)}
\ACNset{level}{متوسط}
\ACNset{prereq}{درس‌های ۱ تا ۳؛ افت مسیر لگاریتمی-فاصله‌ای؛ توزیع گاوسی و تابع $Q$}
\ACNset{gen}{4}

\begin{document}
\ACNcover

% ---------------------------------------------------------------------
\section{اهداف یادگیری}
پس از مطالعهٔ این درس باید بتوانید:
\begin{objectives}
  \item تحرک در حالت بیکار (انتخاب مجدد سلول، ثبت مکان) را از تحرک در حالت متصل (تحویل)
        \textbf{متمایز کنید} و توضیح دهید در هر حالت چه کسی تصمیم می‌گیرد.
  \item ناحیهٔ مکانی، ناحیهٔ مسیریابی، فهرست نواحی ردیابی و ناحیهٔ اعلان شبکهٔ دسترسی را
        \textbf{مقایسه کنید} و اندازهٔ ناحیهٔ ثبت را با یک مدل هزینه \textbf{بهینه کنید}.
  \item تحویل‌ها را \textbf{طبقه‌بندی کنید} (درون‌فرکانسی/میان‌فرکانسی، میان‌فناوری، سخت/نرم) و
        کمیت‌های اندازه‌گیری \acro{RSRP} و \acro{RSRQ} را بیان کنید.
  \item شرط ورود رویداد \acro{A3} در \acro{LTE} را با پسماند، آفست‌ها و زمان تا راه‌اندازی
        \textbf{به کار ببرید} و محل آغاز تحویل را \textbf{محاسبه کنید}.
  \item تحویل مبتنی بر \acro{X2} در \acro{LTE} را پیام‌به‌پیام، شامل ارسال دادهٔ واسطه و تعویض مسیر،
        \textbf{دنبال کنید}.
  \item تحویل‌های دیرهنگام، زودهنگام، به سلول اشتباه و پینگ‌پنگ را \textbf{تشخیص دهید} و آن‌ها را
        به تنظیم پارامترها ربط دهید.
\end{objectives}

% ---------------------------------------------------------------------
\section{دانش پیش‌نیاز}
به مفهوم ناحیهٔ مکانی از درس‌های ۱ و ۲، تحویل نرم \acro{WCDMA} (درس ۳) و مدل افت مسیر
لگاریتمی-فاصله‌ای با محوشدگی سایه‌ای لاگ-نرمال نیاز دارید:
$\PL(d) = \PL(d_0) + 10n\log_{10}(d/d_0) + X_\sigma$، که در آن $X_\sigma$ متغیری گاوسی با میانگین
صفر (بر حسب \lr{dB}) و انحراف معیار $\sigma$ است. تابع دنباله $Q(x) = \Pr\{Z > x\}$ برای متغیر
گاوسی استاندارد $Z$ است.

\paragraph{واژه‌های کلیدی} واژگان این درس همراه با معادل انگلیسی آن‌ها:
\begin{keyterms}
مدیریت تحرک & \lr{Mobility management} & دسترس‌پذیر و متصل نگه داشتن کاربر متحرک. \\
انتخاب مجدد سلول & \lr{Cell reselection} & تغییر سلول اقامت در حالت بیکار به تصمیم \acro{UE}. \\
تحویل & \lr{Handover (HO)} & انتقال اتصال فعال به سلول دیگر به تصمیم شبکه. \\
ناحیهٔ ردیابی & \lr{Tracking area (TA)} & ناحیهٔ ثبت در \lr{LTE/5G}؛ دامنهٔ فراخوانی. \\
توان دریافتی سیگنال مرجع & \lr{RSRP} & میانگین توان یک عنصر منبع حامل سیگنال مرجع سلول. \\
کیفیت دریافتی سیگنال مرجع & \lr{RSRQ} & $N\cdot\mathrm{RSRP}/\mathrm{RSSI}$: کیفیتی آگاه از بار و تداخل. \\
پسماند & \lr{Hysteresis} & حاشیه‌ای که از راه‌اندازی مکرر حول یک آستانه جلوگیری می‌کند. \\
زمان تا راه‌اندازی & \lr{Time-to-trigger (TTT)} & مدتی که شرط باید برقرار بماند تا گزارش ارسال شود. \\
پینگ‌پنگ & \lr{Ping-pong} & تحویل غیرضروری که به‌سرعت با بازگشت دنبال می‌شود. \\
شکست پیوند رادیویی & \lr{Radio link failure (RLF)} & از دست رفتن اتصال که با زمان‌سنج‌ها (مثلاً \lr{T310}) آشکار می‌شود. \\
\end{keyterms}

% ---------------------------------------------------------------------
\section{مقدمه و انگیزه}
تلفن ثابت همیشه به همان پریز وصل است؛ کاربر موبایل نه. مدیریت تحرک به دو پرسش متفاوت پاسخ می‌دهد:
\begin{enumerate}[label=(\alph*)]
  \item \textbf{دسترس‌پذیری} -- وقتی کاربر بیکار است و تماس یا بسته‌ای می‌رسد، شبکه کجا باید
        دنبالش بگردد؟
  \item \textbf{پیوستگی نشست} -- وقتی کاربر فعال است، چگونه اتصال می‌تواند بدون قطع یا وقفه کاربر
        را از سلولی به سلول دیگر دنبال کند؟
\end{enumerate}
هر دو پرسش مصالحه‌ای میان \emph{هزینهٔ سیگنالینگ} و \emph{کارایی} دارند. اگر شبکه سلول دقیق هر
کاربر بیکار را می‌دانست، هر تغییر سلول یک به‌روزرسانی هزینه داشت؛ اگر هیچ نمی‌دانست، هر تماس ورودی
باید همه‌جا فراخوانی می‌شد. اگر تحویل‌ها مشتاقانه آغاز شوند، کاربران میان سلول‌ها جابه‌جا می‌شوند
(پینگ‌پنگ)؛ اگر دیر آغاز شوند، پیوندها قطع می‌شوند. این درس هر دو سو را از نسل دوم تا پنجم
بررسی می‌کند و به هر مصالحه شکلی کمّی می‌دهد.

% ---------------------------------------------------------------------
\section{مفاهیم بنیادی}

\subsection{حالت بیکار در برابر حالت متصل}
در \textbf{حالت بیکار} \acro{UE} اتصال رادیویی ندارد. روی یک سلول \emph{اقامت} می‌کند، به کانال‌های
پخش و فراخوانی آن گوش می‌دهد و خودش تصمیم می‌گیرد چه زمانی به سلول بهتری برود (انتخاب مجدد
سلول). شبکه تنها ناحیهٔ ثبت \acro{UE} را می‌داند. در \textbf{حالت متصل} \acro{UE} اتصال رادیویی
اختصاصی دارد؛ شبکه بر اساس گزارش‌های اندازه‌گیری \acro{UE} سلول خدمت‌دهنده را تعیین می‌کند (تحویل
تحت کنترل شبکه و با کمک موبایل).

\subsection{نواحی ثبت در نسل‌های مختلف}
\begin{itemize}
  \item \textbf{ناحیهٔ مکانی (\acro{LA})} -- حوزهٔ \acro{CS} در \lr{GSM/UMTS}؛ به‌روزرسانی به سوی \acro{VLR}.
  \item \textbf{ناحیهٔ مسیریابی (\acro{RA})} -- حوزهٔ \acro{PS} در \lr{GPRS/UMTS}؛ زیرمجموعه‌ای از یک
        \acro{LA} که به‌روزرسانی آن به سوی \acro{SGSN} است. چون دو حوزه جدا بودند، دو رویهٔ موازی لازم بود.
  \item \textbf{ناحیهٔ ردیابی (\acro{TA})} -- \acro{LTE} و \lr{5G}. \acro{MME} (در \lr{5G}، \acro{AMF}) به
        \acro{UE} یک \emph{فهرست} از \acro{TA}ها می‌دهد؛ \acro{UE} تنها وقتی به‌روزرسانی می‌کند که وارد
        \acro{TA}ای خارج از فهرست شود. فهرست‌ها برای هر \acro{UE} قابل تنظیم‌اند و پینگ‌پنگ
        به‌روزرسانی در یک مرز ثابت را حذف می‌کنند.
  \item \textbf{ناحیهٔ اعلان شبکهٔ دسترسی (\acro{RNA})} -- حالت \lr{RRC\_INACTIVE} در \lr{5G}: هسته
        \acro{UE} را متصل می‌داند، شبکهٔ دسترسی زمینهٔ آن را نگه می‌دارد و در ناحیه‌ای که خودش
        پیکربندی کرده آن را فراخوانی می‌کند.
\end{itemize}

\begin{acnfigure}
  \resizebox{\linewidth}{!}{\acnfig{lecture-04/area-hierarchy}}
  \caption{نواحی ثبت در نسل‌های پیاپی.}
  \label{fig:areas}
\end{acnfigure}

\subsection{کمیت‌های اندازه‌گیری}
\acro{LTE} دو اندازه‌گیری پایه را روی سیگنال‌های مرجع ویژهٔ سلول تعریف می‌کند \cite{ts36300}:
\begin{definitionbox}{\lr{RSRP} و \lr{RSRQ}}
\emph{\acro{RSRP}} میانگین خطی توان عناصر منبعی است که سیگنال‌های مرجع ویژهٔ سلول را در پهنای
باند اندازه‌گیری حمل می‌کنند (بازهٔ تقریبی $-140$ تا $-44$\,\lr{dBm}). \emph{\acro{RSRQ}} برابر است با
\begin{equation}\label{eq:rsrq}
  \RSRQ = \frac{N\cdot\RSRP}{\RSSI},
\end{equation}
که $N$ تعداد بلوک‌های منبع در پهنای باند اندازه‌گیری \acro{RSSI} است و \acro{RSSI} کل توان دریافتی
(سیگنال مطلوب، تداخل و نویز) در آن پهنای باند است.
\end{definitionbox}
\acro{RSRP} پوشش را می‌سنجد؛ \acro{RSRQ} بار و تداخل را نیز منعکس می‌کند، زیرا وقتی همسایه‌ها
مشغول‌اند \acro{RSSI} رشد می‌کند. \acro{NR} از معادل‌های \lr{SS-RSRP} و \lr{SS-RSRQ} روی بلوک‌های
سیگنال همگام‌سازی استفاده می‌کند.

% ---------------------------------------------------------------------
\section{توضیحات فنی تفصیلی}

\subsection{انتخاب مجدد سلول در حالت بیکار}
در \acro{LTE} یک سلول \emph{مناسب} است اگر معیار \lr{S} را برآورده کند \cite{ts36304}:
\begin{equation}\label{eq:scrit}
  S_{\mathrm{rxlev}} = Q_{\mathrm{rxlevmeas}} - (Q_{\mathrm{rxlevmin}} + Q_{\mathrm{rxlevminoffset}})
  - P_{\mathrm{compensation}} > 0 ,
\end{equation}
که $Q_{\mathrm{rxlevmeas}}$ مقدار \acro{RSRP} اندازه‌گیری‌شده، $Q_{\mathrm{rxlevmin}}$ حداقل سطح پخش‌شده و
$P_{\mathrm{compensation}}$ جریمه‌ای برای \acro{UE}هایی است که بیشینهٔ توانشان کمتر از سطح فرض‌شدهٔ
اپراتور است. \acro{UE} در میان سلول‌های مناسب درون‌فرکانسی، سلول خدمت‌دهنده و همسایه‌ها را با
\begin{equation}\label{eq:rcrit}
  R_\mathrm{s} = Q_{\mathrm{meas,s}} + Q_{\mathrm{hyst}}, \qquad
  R_\mathrm{n} = Q_{\mathrm{meas,n}} - Q_{\mathrm{offset}},
\end{equation}
رتبه‌بندی می‌کند و تنها وقتی همسایه‌ای را انتخاب می‌کند که به مدت $T_{\mathrm{reselection}}$ بهترین
رتبه را داشته باشد. انتخاب مجدد میان‌فرکانسی و میان‌فناوری از \emph{اولویت}های مطلق استفاده
می‌کند: \acro{UE} همیشه به لایهٔ با اولویت بالاتر که به اندازهٔ کافی خوب باشد می‌رود و تنها وقتی
لایهٔ فعلی ضعیف شود به لایهٔ با اولویت پایین‌تر می‌رود. اولویت‌ها به اپراتور امکان می‌دهند کاربران
بیکار را مثلاً به سمت یک لایهٔ ظرفیت یا یک حامل \lr{5G} هدایت کند.

\subsection{طبقه‌بندی تحویل}
\begin{itemize}
  \item بر اساس \textbf{فرکانس}: \emph{درون‌فرکانسی} (همان حامل؛ بدون شکاف قابل اندازه‌گیری)،
        \emph{میان‌فرکانسی} (حامل دیگر؛ \acro{UE} به شکاف اندازه‌گیری نیاز دارد مگر گیرندهٔ دوم داشته
        باشد) و \emph{میان‌فناوری} (فناوری دیگر، مثلاً \acro{LTE} به \acro{UMTS}).
  \item بر اساس \textbf{پیوستگی اتصال}: تحویل \emph{سخت} (قطع پیش از اتصال؛ \lr{GSM, LTE, NR}) و
        تحویل \emph{نرم} (اتصال پیش از قطع؛ \acro{WCDMA}، که در آن چند سلول هم‌زمان دادهٔ یکسان ارسال می‌کنند).
  \item بر اساس \textbf{دامنهٔ شبکه}: درون یک کنترل‌گر (درون \acro{BSC}، درون \acro{eNB})، میان
        ایستگاه‌های پایه (تحویل \lr{X2/Xn}) یا با دخالت هسته (میان \acro{MSC}ها در \acro{GSM}، مبتنی بر
        \lr{S1} در \acro{LTE}، مبتنی بر \lr{N2} در \lr{5G}).
\end{itemize}

\paragraph{تحویل در \acro{GSM}} \acro{BTS} فراسو را اندازه می‌گیرد و \acro{MS} شش همسایهٔ قوی‌تر را
هر ۴۸۰ میلی‌ثانیه روی \acro{SACCH} گزارش می‌کند. \acro{BSC} تصمیم می‌گیرد؛ تحویل‌های میان \acro{BSC}ها
\acro{MSC} را درگیر می‌کنند و در تحویل میان \acro{MSC}ها، \acro{MSC} نخست به‌عنوان \emph{لنگر} تماس
باقی می‌ماند \cite{ts23009}.

\paragraph{تحویل نرم در \acro{WCDMA}} \acro{UE} یک \emph{مجموعهٔ فعال} از سلول‌ها نگه می‌دارد؛ رویدادها
سلولی را که در بازهٔ گزارش‌دهی بهترین سلول قرار گیرد اضافه، سلولی را که از آن خارج شود حذف، یا
ضعیف‌ترین عضو را جایگزین می‌کنند. هزینهٔ آن توان فروسوی اضافی و ظرفیت \acro{Iub} است؛ معمولاً ۳۰ تا
۴۰ درصد کاربران در تحویل نرم‌اند.

\paragraph{تحویل در \acro{LTE}} \acro{LTE} \acro{RNC} و تحویل نرم را حذف کرد. \acro{eNB} به‌تنهایی و بر
پایهٔ \emph{رویدادهای} اندازه‌گیری که از طریق \acro{RRC} پیکربندی می‌شوند تصمیم می‌گیرد \cite{ts36331}:
\begin{table}[htbp]
\centering\small
\caption{رویدادهای اندازه‌گیری درون‌فناوری \acro{LTE} (آستانه‌ها و آفست‌ها را \acro{RRC} پیکربندی می‌کند).}\label{tab:events}
\rowcolors{2}{acnmist}{white}
\begin{tabularx}{\linewidth}{@{}l Y Y@{}}
\toprule
\rowcolor{white}رویداد & شرط ورود (ساده‌شده) & کاربرد رایج \\
\midrule
\lr{A1} & سلول خدمت‌دهنده از یک آستانه بهتر می‌شود & توقف اندازه‌گیری میان‌فرکانسی (حذف شکاف‌ها) \\
\lr{A2} & سلول خدمت‌دهنده از یک آستانه بدتر می‌شود & آغاز اندازه‌گیری میان‌فرکانسی/میان‌فناوری \\
\lr{A3} & همسایه به اندازهٔ آفست از سلول خدمت‌دهنده (\acro{PCell}) بهتر می‌شود & تحویل درون‌فرکانسی \\
\lr{A4} & همسایه از یک آستانه بهتر می‌شود & تحویل مبتنی بر بار یا لایهٔ پوشش \\
\lr{A5} & خدمت‌دهنده بدتر از آستانهٔ ۱ و همسایه بهتر از آستانهٔ ۲ & تحویل میان‌فرکانسی \\
\lr{A6} & همسایه به اندازهٔ آفست از یک سلول ثانویه بهتر می‌شود & مدیریت تجمیع حامل \\
\lr{B1/B2} & همسایهٔ فناوری دیگر بهتر از آستانه (و در \lr{B2} خدمت‌دهنده بدتر) & \lr{LTE} به \lr{UMTS/GSM}، \lr{LTE} به \lr{NR} \\
\bottomrule
\end{tabularx}
\end{table}

\subsection{تحویل در \lr{5G NR}}
\acro{NR} اصول \acro{LTE} را حفظ می‌کند (\lr{eNB$\to$gNB}، \lr{X2$\to$Xn}، \lr{S1$\to$N2}) و ویژگی‌هایی
برای کاهش وقفه و شکست‌های تحویل سخت می‌افزاید \cite{ts38300}:
\begin{itemize}
  \item \textbf{تحویل شرطی (\acro{CHO})}: شبکه یک یا چند سلول مقصد را از پیش آماده می‌کند و فرمان تحویل
        را همراه با یک شرط اجرا به \acro{UE} می‌فرستد؛ \acro{UE} وقتی شرط برقرار شد آن را اجرا می‌کند و
        از این رقابت جلوگیری می‌شود که فرمان روی پیوندی در حال شکست برسد.
  \item \textbf{پشتهٔ پروتکل فعال دوگانه (\acro{DAPS})}: \acro{UE} در حین اتصال به مقصد همچنان از مبدأ
        دریافت می‌کند و وقفه به صفر نزدیک می‌شود.
  \item \textbf{مدیریت پرتو} درون سلول: بسیاری از تغییرات پرتو در لایه‌های ۱ و ۲ و بدون تحویل
        \acro{RRC} انجام می‌شوند.
\end{itemize}

\subsection{شکست‌های تحویل}
\acro{UE} وقتی \emph{شکست پیوند رادیویی} (\acro{RLF}) اعلام می‌کند که پس از $N310$ نشانهٔ پیاپی
خارج از همگامی، زمان‌سنج $T310$ بدون بازیابی به پایان برسد؛ سپس برقراری مجدد \acro{RRC} را امتحان
می‌کند. شبکه‌های خودسازمان‌ده شکست‌ها را برای تنظیم پارامترها طبقه‌بندی می‌کنند (بهینه‌سازی
استحکام تحرک، \acro{MRO}):
\begin{itemize}
  \item \textbf{دیرهنگام} -- \acro{RLF} در سلول مبدأ پیش از (یا در حین) تحویل؛ \acro{UE} در مقصد برقراری
        مجدد می‌کند. راه‌حل: راه‌اندازی زودتر (آفست یا \acro{TTT} کمتر).
  \item \textbf{زودهنگام} -- \acro{RLF} کوتاه‌زمانی پس از تحویل موفق؛ \acro{UE} در \emph{مبدأ} برقراری
        مجدد می‌کند. راه‌حل: راه‌اندازی دیرتر.
  \item \textbf{سلول اشتباه} -- \acro{RLF} کوتاه‌زمانی پس از تحویل و برقراری مجدد در سلول سوم.
        راه‌حل: تنظیم آفست‌های ویژهٔ سلول دو همسایه.
  \item \textbf{پینگ‌پنگ} -- تحویل موفقی که در زمان کوتاهی با بازگشت به مبدأ دنبال می‌شود؛ سیگنالینگ
        را هدر می‌دهد و وقفه ایجاد می‌کند. راه‌حل: پسماند یا \acro{TTT} بیشتر.
\end{itemize}
تحویل دیرهنگام و پینگ‌پنگ پارامترها را در جهت‌های مخالف می‌کشند: این دوراهی اصلی تنظیم تحویل است.

% ---------------------------------------------------------------------
\section{مدل‌های ریاضی}

\subsection{شرط \lr{A3} و نقطهٔ تحویل}\label{sec:a3}
\paragraph{مسئلهٔ مهندسی} تحویل واقعاً در کجای جاده‌ای میان دو سلول رخ می‌دهد و پارامترها این نقطه
را چگونه جابه‌جا می‌کنند؟
\paragraph{شرط ورود \lr{A3}} در \acro{LTE} \cite{ts36331} این شرط، با همهٔ کمیت‌ها بر حسب \lr{dB} یا \lr{dBm}، عبارت است از
\begin{equation}\label{eq:a3}
  M_\mathrm{n} + O_{\mathrm{fn}} + O_{\mathrm{cn}} - \Hys \;>\; M_\mathrm{p} + O_{\mathrm{fp}} + O_{\mathrm{cp}} + \Off ,
\end{equation}
که $M_\mathrm{n}$ و $M_\mathrm{p}$ اندازه‌گیری‌های (فیلترشدهٔ لایهٔ ۳) همسایه و سلول خدمت‌دهندهٔ اصلی،
$O_{\mathrm{f}\cdot}$ آفست‌های ویژهٔ فرکانس، $O_{\mathrm{c}\cdot}$ آفست‌های ویژهٔ سلول، $\Hys$ پسماند و
$\Off$ آفست رویداد هستند. شرط باید به مدت $\TTT$ پیوسته برقرار باشد تا \acro{UE} گزارش دهد. با
آفست‌های فرکانس و سلول صفر، به $M_\mathrm{n} - M_\mathrm{p} > H$ با $H = \Off + \Hys$ ساده می‌شود.

\paragraph{مدل سامانه} دو \acro{eNB} در فاصلهٔ $D$ با توان ارسال برابر؛ \acro{UE} روی خط واصل آن‌ها
حرکت می‌کند؛ افت مسیر قطعی با نمای $n$ (بدون سایه). در فاصلهٔ $x$ از \acro{eNB} خدمت‌دهنده،
\begin{equation}
  M_\mathrm{n} - M_\mathrm{p} = 10n\log_{10}\frac{x}{D-x}.
\end{equation}
\paragraph{استخراج} با قرار دادن $10n\log_{10}\bigl(x/(D-x)\bigr) = H$ و حل آن،
\begin{equation}\label{eq:hopoint}
  x_{\mathrm{HO}} = D\,\frac{r}{1+r}, \qquad r = 10^{H/(10n)} .
\end{equation}
\acro{UE} سپس پیش از گزارش‌دادن مسافت $v\,\TTT$ دیگر را طی می‌کند.
\paragraph{تفسیر} برای $D = \qty{1000}{\metre}$، $n=3.5$ و $H = \qty{3}{\dB}$: $r=1.218$ و
$x_{\mathrm{HO}} = \qty{549}{\metre}$، یعنی ۴۹ متر پس از مرز سلول؛ $H = \qty{6}{\dB}$ آن را به ۵۹۷ متر
می‌برد. با سرعت \qty{30}{\metre\per\second}، \acro{TTT} برابر ۳۲۰ میلی‌ثانیه ۹٫۶ متر دیگر می‌افزاید
(\cref{fig:a3}). بهای هر دسی‌بل حاشیه با کیفیت سیگنال روی سلول قدیمی پرداخت می‌شود، جایی که پیوند
از همه ضعیف‌تر است -- منشأ شکست‌های دیرهنگام در سرعت بالا.

\begin{acnfigure}
  \acnfig{lecture-04/a3-trigger}
  \caption{راه‌اندازی \lr{A3}: همسایه باید به مدت $\TTT = 320$\,\lr{ms} به اندازهٔ
  $\Off + \Hys = 3$\,\lr{dB} (نقطه‌چین) از سلول خدمت‌دهنده بهتر باشد تا گزارش اندازه‌گیری ارسال شود.}
  \label{fig:a3}
\end{acnfigure}

\subsection{سایه و پینگ‌پنگ}\label{sec:pingpong}
\paragraph{مدل} سایهٔ لاگ-نرمال مستقل $X_\mathrm{s}, X_\mathrm{n}\sim\mathcal{N}(0,\sigma^2)$ (\lr{dB}) را
به هر دو سلول اضافه کنید. در مرز سلول اختلاف میانگین صفر است و
$M_\mathrm{n} - M_\mathrm{p} \sim \mathcal{N}(0, 2\sigma^2)$.
\paragraph{نتیجه} احتمال اینکه یک نمونه در مرز از حاشیهٔ $H$ بیشتر شود برابر است با
\begin{equation}\label{eq:pingpong}
  P_{\mathrm{trig}} = Q\!\left(\frac{H}{\sqrt{2}\,\sigma}\right).
\end{equation}
با $\sigma = \qty{8}{\dB}$ و $H = \qty{3}{\dB}$، $P_{\mathrm{trig}} = Q(0.265) = 0.40$؛ حتی با
$H=\qty{6}{\dB}$ هنوز $0.30$ است. پسماند به‌تنهایی نمی‌تواند راه‌اندازی‌های کاذب ناشی از محوشدگی
بزرگ‌مقیاس را سرکوب کند؛ \acro{TTT} و فیلتر لایهٔ ۳ لازم‌اند. فیلتر لایهٔ ۳
\begin{equation}\label{eq:l3filter}
  F_k = (1-a)\,F_{k-1} + a\,M_k, \qquad a = 2^{-K/4},
\end{equation}
با ضریب فیلتر $K$، واریانس نوسان‌های \emph{مستقل} (محوشدگی سریع، خطای اندازه‌گیری) را با ضریب
$a/(2-a)$ کاهش می‌دهد -- یک‌سوم برای $K=4$ -- اما نمی‌تواند سایه را که در طول ده‌ها متر همبسته است حذف
کند. این مدل از یک جهت بدبینانه است (همبستگی سایهٔ دو سلول را نادیده می‌گیرد) و از جهت دیگر
خوش‌بینانه (تنها یک نقطه را بررسی می‌کند)؛ گام بعدی معمول، شبیه‌سازی در طول مسیر حرکت است.

\subsection{اندازه‌گذاری ناحیهٔ ثبت}\label{sec:lasize}
\paragraph{مسئله} یک ناحیهٔ مکانی یا ردیابی باید چند سلول داشته باشد؟
\paragraph{مدل و فرض‌ها} سلول‌های مربعی به ضلع $a$؛ ناحیهٔ ثبتی با $N$ سلول مربعی با ضلع $a\sqrt{N}$ و
محیط $L = 4a\sqrt{N}$ می‌سازد. کاربران به‌طور یکنواخت توزیع شده‌اند و در جهت‌های تصادفی یکنواخت با
سرعت میانگین $v$ حرکت می‌کنند (مدل \emph{جریان سیال} \cite{akyildiz1999mobility}): نرخ عبور کاربران
ناحیه‌ای به مساحت $S$ از محیط $L$ آن $\rho vL/\pi$ است که $\rho$ چگالی کاربران است. هر کاربر $\lambda_\mathrm{c}$
تماس در واحد زمان دریافت می‌کند که هر یک در همهٔ $N$ سلول فراخوانی می‌شود. هر به‌روزرسانی $c_\mathrm{u}$
و فراخوانی در هر سلول $c_\mathrm{p}$ هزینه دارد.
\paragraph{استخراج} برای هر کاربر نرخ به‌روزرسانی $\rho vL/(\pi\rho S) = vL/(\pi S)$ با $S = Na^2$ است؛ پس
\begin{equation}\label{eq:cost}
  C(N) = c_\mathrm{u}\,\frac{4v}{\pi a\sqrt{N}} + c_\mathrm{p}\,\lambda_\mathrm{c}\,N .
\end{equation}
جملهٔ اول با $N$ کاهش و جملهٔ دوم افزایش می‌یابد. با قرار دادن $\mathrm{d}C/\mathrm{d}N = 0$،
\begin{equation}\label{eq:nopt}
  -\frac{2c_\mathrm{u}v}{\pi a}N^{-3/2} + c_\mathrm{p}\lambda_\mathrm{c} = 0
  \quad\Longrightarrow\quad
  N^\ast = \left(\frac{2\,c_\mathrm{u}\,v}{\pi\,a\,c_\mathrm{p}\,\lambda_\mathrm{c}}\right)^{2/3}.
\end{equation}
\paragraph{تفسیر} کاربران سریع ($v$ بزرگ) نواحی بزرگ را ترجیح می‌دهند؛ کاربرانی که تماس زیادی
دریافت می‌کنند ($\lambda_\mathrm{c}$ بزرگ) نواحی کوچک را. چون نقطهٔ بهینه به کاربر وابسته است، فهرست‌های
ناحیهٔ ردیابی ویژهٔ هر \acro{UE} در \acro{LTE} می‌توانند کاربر به کاربر به آن نزدیک شوند، کاری که یک
طرح ثابت \acro{LA} نمی‌تواند.

\begin{acnfigure}
  \acnfig{lecture-04/la-cost}
  \caption{هزینهٔ سیگنالینگ به‌روزرسانی، فراخوانی و مجموع در \cref{eq:cost} برای
  $c_\mathrm{u} = 1$، $c_\mathrm{p} = 0.1$، $v = \qty{10}{\metre\per\second}$، $a = \qty{500}{\metre}$ و
  $\lambda_\mathrm{c} = 2$ تماس در ساعت (\cref{ex:la}).}
  \label{fig:lacost}
\end{acnfigure}

% ---------------------------------------------------------------------
\section{معماری شبکه و نمودارها}
معماری تحویل در \acro{LTE} تخت است: \acro{eNB}های مبدأ و مقصد مستقیماً از طریق \acro{X2} مذاکره
می‌کنند و \acro{MME} و \lr{S-GW} تنها در پایان برای جابه‌جایی مسیر فروسو درگیر می‌شوند (\emph{تعویض
مسیر}). اگر \acro{X2} وجود نداشته باشد یا \acro{MME} یا \lr{S-GW} باید عوض شوند، همان گام‌ها از طریق
\acro{MME} روی \lr{S1} اجرا می‌شوند (\emph{تحویل مبتنی بر \lr{S1}}). در \lr{5G} نقش‌ها به \acro{gNB}،
\acro{Xn}، \lr{AMF/SMF} و \acro{UPF} نگاشت می‌شوند (درس ۲۰). \acro{GSM} و \acro{UMTS} کنترل‌گرهای خود
(\lr{BSC/RNC}) و در مرز کنترل‌گرها \lr{MSC/SGSN} را درگیر می‌کنند.

% ---------------------------------------------------------------------
\section{رویه‌های گام‌به‌گام}

\begin{procedurebox}{تحویل مبتنی بر \lr{X2} در \lr{LTE}}\label{proc:x2}
\begin{steps}
  \item \textbf{آغاز آماده‌سازی.} \acro{UE} پس از برقراری شرط \lr{A3} به مدت \acro{TTT} یک
        \msg{Measurement Report} می‌فرستد. \acro{eNB} مبدأ تصمیم به تحویل می‌گیرد.
  \item \textbf{آماده‌سازی.} مبدأ پیام \lr{X2AP} \msg{Handover Request} را همراه با زمینهٔ \acro{UE}
        (حامل‌ها، امنیت، پیکربندی \acro{RRC}) می‌فرستد. مقصد کنترل پذیرش را انجام می‌دهد، منابع
        (از جمله یک پیش‌نوای اختصاصی \acro{RACH}) رزرو می‌کند و با \msg{Handover Request Acknowledge}
        که پیام \acro{RRC} برای \acro{UE} را در بر دارد پاسخ می‌دهد.
  \item \textbf{اجرا.} مبدأ پیام \msg{RRC Connection Reconfiguration} را همراه با اطلاعات کنترل تحرک
        می‌فرستد، زمان‌بندی \acro{UE} را متوقف می‌کند، \msg{SN Status Transfer} را می‌فرستد و بسته‌های
        فروسوی بافرشده و در راه را از طریق \acro{X2} به مقصد ارسال می‌کند.
  \item \textbf{دسترسی.} \acro{UE} با مقصد همگام می‌شود، دسترسی تصادفی بدون رقابت (درس ۸) انجام
        می‌دهد، پیش‌فرستی زمانی و یک مجوز فراسو دریافت می‌کند و \msg{RRC Connection Reconfiguration Complete}
        را می‌فرستد. وقفه پایان می‌یابد.
  \item \textbf{تکمیل.} مقصد با \lr{S1AP} \msg{Path Switch Request} از \acro{MME} می‌خواهد مسیر را جابه‌جا
        کند؛ \acro{MME} پیام \msg{Modify Bearer Request} را به \lr{S-GW} می‌فرستد که تونل فروسو را عوض
        می‌کند و بسته‌های \emph{نشانگر پایان} را روی مسیر قدیمی می‌فرستد تا مقصد بداند دادهٔ ارسال‌شده
        تمام شده است.
  \item \textbf{آزادسازی.} پس از \msg{Path Switch Request Acknowledge}، مقصد \msg{UE Context Release} را
        می‌فرستد و مبدأ منابعش را آزاد می‌کند.
\end{steps}
\end{procedurebox}

\begin{acnfigure}
  \resizebox{\linewidth}{!}{\acnfig{lecture-04/x2-handover}}
  \caption{تحویل مبتنی بر \lr{X2} در \acro{LTE} (\cref{proc:x2}). پیکان‌های فیروزه‌ای: ارسال دادهٔ
  صفحهٔ کاربر و نشانگر پایان.}
  \label{fig:x2}
\end{acnfigure}

وقفهٔ صفحهٔ کاربر از لحظه‌ای که \acro{UE} مبدأ را ترک می‌کند تا وقتی که بتواند در مقصد زمان‌بندی
شود ادامه دارد؛ در \acro{LTE} از مرتبهٔ چند ده میلی‌ثانیه است و همان کمیتی است که تحویل بدون
\acro{RACH}، اتصال پیش از قطع و \acro{DAPS} کاهش می‌دهند.

% ---------------------------------------------------------------------
\section{مثال‌های حل‌شده}

\begin{examplebox}{نقطهٔ تحویل و سرعت}\label{ex:hopoint}
\textbf{مسئله.} فاصلهٔ میان سایت‌ها ۱۰۰۰ متر، $n=3.5$، $\Off+\Hys = 3$\,\lr{dB} و
$\TTT = 320$\,\lr{ms}. گزارش اندازه‌گیری برای یک عابر پیاده (\qty{1.5}{\metre\per\second}) و یک خودرو در
بزرگراه (\qty{120}{\kilo\metre\per\hour}) در کجا ارسال می‌شود؟

\textbf{حل.} از \cref{eq:hopoint}، $x_{\mathrm{HO}} = 549$\,\lr{m}. عابر $1.5\times0.32 = 0.5$ متر
می‌افزاید؛ خودرو $33.3\times0.32 = 10.7$ متر و در حدود \qty{560}{\metre} گزارش می‌دهد.

\textbf{تفسیر.} سپس آماده‌سازی و اجرا چند ده میلی‌ثانیه دیگر طول می‌کشد و خودرو در هر ۳۰ میلی‌ثانیه
۱ تا ۳ متر دیگر جابه‌جا می‌شود. به همین دلیل به سلول‌های بزرگراه اغلب آفست یا \acro{TTT} کمتری داده
می‌شود و مقیاس‌گذاری \acro{TTT} وابسته به سرعت بخشی از استاندارد است.
\end{examplebox}

\begin{examplebox}{\lr{RSRQ} تحت بار}\label{ex:rsrq}
\textbf{مسئله.} یک \acro{UE} مقدار $\RSRP=-90$\,\lr{dBm} را روی $N=50$ بلوک منبع اندازه می‌گیرد. \acro{RSRQ}
را وقتی \acro{RSSI} برابر $-65$\,\lr{dBm} (همسایه‌های کم‌بار) و $-60$\,\lr{dBm} (پربار) است حساب کنید.

\textbf{حل.} $\RSRQ = 10\log_{10}50 + \RSRP - \RSSI = 17.0 - 90 + 65 = \qty{-8.0}{\dB}$ و
$17.0 - 90 + 60 = \qty{-13.0}{\dB}$.

\textbf{تفسیر.} پوشش یکسان، کیفیت ۵ دسی‌بل بدتر: \acro{RSRQ} بار همسایه‌ها را می‌بیند و \acro{RSRP}
نمی‌بیند. استفاده از \acro{RSRQ} برای \lr{A3} بار را پخش می‌کند، به بهای گزارش‌های پرنوسان‌تر.
\end{examplebox}

\begin{examplebox}{اندازهٔ ناحیهٔ ردیابی}\label{ex:la}
\textbf{مسئله.} با پارامترهای \cref{fig:lacost}، تعداد بهینهٔ سلول در هر ناحیه و هزینه به ازای کاربر
در ساعت را بیابید.

\textbf{حل.} $\lambda_\mathrm{c} = 2/3600\,\mathrm{s}^{-1}$. طبق \cref{eq:nopt}:
$N^\ast = \bigl(2\times1\times10/(\pi\times500\times0.1\times2/3600)\bigr)^{2/3} = (229.2)^{2/3} = 37.5$.
یک ناحیهٔ $6\times6$ ($N=36$) هزینهٔ $C = 3600\,[4\times10/(\pi\times500\times6) + 0.1\times(2/3600)\times36]
= 15.28 + 7.20 = 22.5$ واحد به ازای کاربر-ساعت دارد، در برابر ۹۱٫۹ برای $N=1$ و ۲۹٫۲ برای $N=100$.

\textbf{تفسیر.} نقطهٔ بهینه تخت است: هر مقداری میان حدود ۲۰ تا ۶۰ سلول کمتر از ۱۰٪ بیش از کمینه
هزینه دارد. بنابراین اپراتورها می‌توانند مرزهای ناحیه را به دلایل عملی (در امتداد خطوط کم‌ترافیک،
دور از جاده‌های شلوغ) بدون زیان زیاد انتخاب کنند.
\end{examplebox}

% ---------------------------------------------------------------------
\section{مقایسهٔ فناوری‌ها}
\begin{table}[htbp]
\centering\small
\caption{تحرک در حالت متصل از نسل دوم تا پنجم.}\label{tab:compare}
\rowcolors{2}{acnmist}{white}
\begin{tabularx}{\linewidth}{@{}l Y Y Y Y@{}}
\toprule
\rowcolor{white} & \acro{GSM} & \acro{UMTS} (\acro{WCDMA}) & \acro{LTE} & \acro{NR} (\lr{5G}) \\
\midrule
نوع               & سخت                  & نرم (هم‌فرکانس) + سخت       & سخت                     & سخت؛ گزینه‌های \lr{CHO} و \lr{DAPS} \\
تصمیم‌گیرنده      & \acro{BSC}           & \acro{RNC}                  & \acro{eNB} مبدأ          & \acro{gNB} مبدأ \\
اندازه‌گیری‌ها     & \lr{RXLEV, RXQUAL}   & \lr{CPICH RSCP}، $E_\mathrm{c}/N_0$ & \lr{RSRP, RSRQ, SINR} & \lr{SS-RSRP, SS-RSRQ, SS-SINR, CSI-RS} \\
پیوند میان گره‌ها & \lr{A} (از طریق \lr{MSC}) & \lr{Iur}               & \lr{X2} (یا \lr{S1})     & \lr{Xn} (یا \lr{N2}) \\
دخالت هسته        & لنگر \acro{MSC}      & جابه‌جایی \acro{SRNC}       & تعویض مسیر (\lr{MME, S-GW}) & تعویض مسیر (\lr{AMF, SMF, UPF}) \\
ناحیهٔ بیکار      & \acro{LA} (\acro{RA} برای \acro{PS}) & \lr{LA + RA} & فهرست \acro{TA}         & فهرست \acro{TA} + \acro{RNA} \\
\bottomrule
\end{tabularx}
\end{table}

\paragraph{سخت در برابر نرم} تحویل نرم بهرهٔ تنوع و بدون وقفه بودن را می‌دهد، اما به همسایه‌های
هم‌فرکانس، یک نقطهٔ ترکیب مرکزی (\acro{RNC}) و منابع اضافی نیاز دارد. \acro{LTE} معماری تخت بدون گرهٔ
مرکزی در شبکهٔ دسترسی را برگزید و تحویل نرم ناممکن شد؛ این را با تحویل سخت سریع، ارسال دادهٔ واسطه
و بعدها \acro{DAPS} جبران کرد. این مصالحه -- \emph{ترکیب متمرکز در برابر تصمیم‌های تخت و سریع} -- در
\lr{5G} با تقسیم \lr{CU/DU} (درس ۱۴) بازمی‌گردد.

% ---------------------------------------------------------------------
\section{نکته‌های کلیدی}
\begin{takeawaybox}
\begin{itemize}
  \item تحرک بیکار را \acro{UE} کنترل می‌کند (انتخاب مجدد، به‌روزرسانی ناحیه)؛ تحرک متصل تحت کنترل شبکه
        و با کمک موبایل است (گزارش اندازه‌گیری، تحویل).
  \item اندازهٔ ناحیه هزینهٔ به‌روزرسانی را با هزینهٔ فراخوانی معاوضه می‌کند:
        $N^\ast = (2c_\mathrm{u}v/(\pi a c_\mathrm{p}\lambda_\mathrm{c}))^{2/3}$؛ فهرست‌های \acro{TA} به شبکه امکان
        می‌دهند آن را برای هر \acro{UE} تنظیم کند.
  \item \acro{RSRP} پوشش را می‌سنجد و $\RSRQ = N\,\RSRP/\RSSI$ بار و تداخل را نیز در نظر می‌گیرد.
  \item \lr{A3}: $M_\mathrm{n} - M_\mathrm{p} > \Off + \Hys$ (به‌علاوهٔ آفست‌ها) به مدت $\TTT$؛ نقطهٔ تحویل به
        $x = Dr/(1+r)$ با $r = 10^{H/(10n)}$ منتقل می‌شود.
  \item با سایهٔ \lr{8\,dB}، حاشیهٔ \lr{3\,dB} در ۴۰٪ موارد در مرز به‌طور تصادفی عبور می‌شود: \acro{TTT} و
        فیلترکردن ضروری‌اند.
  \item تحویل \acro{X2}: گزارش $\to$ درخواست/تأیید $\to$ فرمان $\to$ ارسال داده $\to$ \acro{RACH}
        $\to$ تکمیل $\to$ تعویض مسیر $\to$ آزادسازی.
  \item تحویل دیرهنگام و پینگ‌پنگ پارامترها را در جهت‌های مخالف می‌کشند؛ \acro{MRO}، \acro{CHO} و
        \acro{DAPS} به این دوراهی می‌پردازند.
\end{itemize}
\end{takeawaybox}

% ---------------------------------------------------------------------
\section{تمرین‌ها}

\begin{exercisebox}{1}{چه کسی تصمیم می‌گیرد؟}\label{exr:who}
برای هر یک از موارد زیر بگویید \acro{UE} تصمیم می‌گیرد یا شبکه و بر پایهٔ چه اطلاعاتی: (الف) انتخاب
مجدد سلول در \acro{LTE}؛ (ب) تحویل در \acro{LTE}؛ (ج) به‌روزرسانی مجموعهٔ فعال در \acro{WCDMA}؛
(د) اجرای تحویل شرطی در \lr{5G}.
\end{exercisebox}

\begin{exercisebox}{1}{دو رویهٔ به‌روزرسانی}\label{exr:lara}
توضیح دهید چرا یک گوشی \acro{GSM} با قابلیت \acro{GPRS} ممکن بود در یک مرز هم به‌روزرسانی ناحیهٔ
مکانی و هم به‌روزرسانی ناحیهٔ مسیریابی انجام دهد، و \acro{LTE} چگونه از این دوگانگی پرهیز کرد.
\end{exercisebox}

\begin{exercisebox}{2}{آفست‌ها}\label{exr:offsets}
در \cref{ex:hopoint} اپراتور برای سبک کردن بار سلول خدمت‌دهنده، آفست ویژهٔ سلول
$O_\mathrm{cn} = +2$\,\lr{dB} را برای همسایه تنظیم می‌کند. اکنون شرط \lr{A3} در کجا برقرار می‌شود؟
احتمال کدام نوع شکست بیشتر می‌شود؟
\end{exercisebox}

\begin{exercisebox}{2}{فیلترکردن}\label{exr:filter}
اندازه‌گیری‌ها هر ۴۰ میلی‌ثانیه با خطاهای مستقل با انحراف معیار ۴ دسی‌بل انجام می‌شوند. با استفاده از
\cref{eq:l3filter}، انحراف معیار پس از فیلتر را برای $K=4$ و $K=8$ و ثابت زمانی فیلتر (تعداد نمونه‌ها
تا افت پاسخ ضربه به $1/\e$) را محاسبه کنید.
\end{exercisebox}

\begin{exercisebox}{3}{ناحیهٔ بهینه با دو رده}\label{exr:twoclass}
یک ناحیهٔ ردیابی به ۷۰٪ عابر پیاده ($v=1.5$\,\lr{m/s}، ۱ تماس در ساعت) و ۳۰٪ خودرو ($v=15$\,\lr{m/s}،
۳ تماس در ساعت) خدمت می‌دهد. با \cref{eq:cost} و بقیهٔ پارامترهای \cref{ex:la}، یک مقدار $N$ که
میانگین هزینه را کمینه کند و $N$ بهینه را در صورتی که هر رده فهرست \acro{TA} خودش را داشته باشد بیابید.
تنظیم ویژهٔ هر رده چقدر صرفه‌جویی می‌کند؟
\end{exercisebox}

\begin{exercisebox}{3}{تحویل دیرهنگام}\label{exr:late}
فرض کنید اگر \acro{RSRP} خدمت‌دهنده پیش از دریافت فرمان تحویل ۸ دسی‌بل پایین‌تر از همسایه برود
\acro{RLF} رخ می‌دهد، و آماده‌سازی و تحویل فرمان ۸۰ میلی‌ثانیه پس از گزارش طول می‌کشد. برای هندسهٔ
\cref{ex:hopoint} بدون سایه، بیشینهٔ سرعتی را بیابید که با $H = 3$\,\lr{dB} و $\TTT = 320$\,\lr{ms} هیچ
شکست دیرهنگامی رخ ندهد. بحث کنید برای خط قطار پرسرعت کدام پارامتر را تغییر می‌دهید.
\end{exercisebox}

\subsection*{پاسخ‌ها}
\addcontentsline{toc}{subsection}{پاسخ‌ها}

\begin{solutionbox}{\ref{exr:who}}
(الف) \acro{UE}، با پارامترهای پخش‌شده ($Q_{\mathrm{rxlevmin}}$، $Q_{\mathrm{hyst}}$، اولویت‌ها) و
اندازه‌گیری‌های خودش. (ب) \acro{eNB} مبدأ، با گزارش‌های اندازه‌گیری \acro{UE} و دانش خود از بار و پذیرش.
(ج) \acro{RNC}، با گزارش رویدادهای \acro{UE} (\lr{1A/1B/1C}). (د) \acro{UE} اجرا می‌کند، اما تنها به سوی
مقصدهایی که شبکه آماده کرده و وقتی شرط پیکربندی‌شده توسط شبکه برقرار شود -- تصمیمی مشترک.
\end{solutionbox}

\begin{solutionbox}{\ref{exr:lara}}
\acro{LA} و \acro{RA} به دو حوزهٔ هستهٔ جدا (\lr{MSC/VLR} و \acro{SGSN}) با رویه‌های جدا تعلق داشتند.
مگر اینکه واسط اختیاری \acro{Gs} به \acro{SGSN} اجازه می‌داد از طرف گوشی \acro{MSC} را به‌روز کند
(به‌روزرسانی ترکیبی)، گوشی دو بار سیگنالینگ می‌کرد. \acro{LTE} یک هستهٔ بسته‌ای واحد با یک موجودیت
تحرک (\acro{MME}) و یک مفهوم ناحیه (\acro{TA}) دارد؛ ثبت مربوط به \acro{CS}، هرجا برای بازگشت به \acro{CS}
لازم باشد، توسط \acro{MME} روی \lr{SGs} انجام می‌شود.
\end{solutionbox}

\begin{solutionbox}{\ref{exr:offsets}}
شرط به $M_\mathrm{n} + 2 - M_\mathrm{p} > 3$ تبدیل می‌شود، یعنی $H_{\mathrm{eff}} = 1$\,\lr{dB}:
$r = 10^{1/35} = 1.068$ و $x_{\mathrm{HO}} = 1000\times1.068/2.068 = 516$\,\lr{m}. تحویل‌ها زودتر رخ می‌دهند؛
احتمال پینگ‌پنگ و شکست زودهنگام بیشتر و احتمال شکست دیرهنگام کمتر می‌شود.
\end{solutionbox}

\begin{solutionbox}{\ref{exr:filter}}
ضریب واریانس $a/(2-a)$ است. $K=4$: $a=0.5$، ضریب $1/3$، $\sigma_F = 4/\sqrt3 = 2.31$\,\lr{dB}.
$K=8$: $a = 0.25$، ضریب $1/7$، $\sigma_F = 1.51$\,\lr{dB}. پاسخ ضربهٔ $(1-a)^k$ پس از
$k = -1/\ln(1-a)$ نمونه به $1/\e$ می‌رسد: ۱٫۴۴ نمونه (۵۸ میلی‌ثانیه) برای $K=4$ و ۳٫۴۸ نمونه (۱۳۹
میلی‌ثانیه) برای $K=8$. هموارسازی بهتر به بهای کندتر شدن واکنش است.
\end{solutionbox}

\begin{solutionbox}{\ref{exr:twoclass}}
هزینهٔ هر کاربر $C_i(N) = 4c_\mathrm{u}v_i/(\pi a\sqrt N) + c_\mathrm{p}\lambda_iN$ است. میانگین همین شکل را با
$\bar v = 0.7\times1.5+0.3\times15 = 5.55$\,\lr{m/s} و $\bar\lambda = 0.7\times1+0.3\times3 = 1.6$ تماس در ساعت
دارد؛ پس $N^\ast = (2\times5.55/(\pi\times500\times0.1\times1.6/3600))^{2/3} = 29.3$ با هزینهٔ حدود
۱۴٫۱ واحد به ازای کاربر-ساعت. به‌طور جداگانه: عابران $N^\ast = 16.8$ (هزینهٔ ۵٫۰)، خودروها
$N^\ast = 37.5$ (هزینهٔ ۳۳٫۷)، میانگین $0.7\times5.0+0.3\times33.7 = 13.6$. تنظیم ویژهٔ هر رده در اینجا تنها
حدود ۳٪ صرفه‌جویی می‌کند، زیرا منحنی هزینه نزدیک کمینه تخت است؛ وقتی رده‌ها بسیار متفاوت‌تر باشند
(مثلاً دستگاه‌های ثابت اینترنت اشیا در برابر قطارها) سود بیشتری دارد.
\end{solutionbox}

\begin{solutionbox}{\ref{exr:late}}
نقطهٔ \acro{RLF} جایی است که $M_\mathrm{n}-M_\mathrm{p} = 8$\,\lr{dB}: $r = 10^{8/35} = 1.693$ و
$x_{\mathrm{RLF}} = 1000\times1.693/2.693 = 629$\,\lr{m}. گزارش در ۵۴۹ متر راه‌اندازی می‌شود و فرمان پس از
$\TTT + 80$\,\lr{ms} $= 0.40$\,\lr{s} می‌رسد. \acro{UE} نباید در این مدت بیش از $629-549 = 80$ متر حرکت کند:
$v_{\max} = 80/0.40 = 200$\,\lr{m/s}. بدون سایه حاشیه بسیار زیاد است؛ با سایه، محوشدگی سلول خدمت‌دهنده
می‌تواند این فاصله را به‌شدت کاهش دهد. برای خط قطار می‌توان \acro{TTT} را کاهش داد (مقیاس‌گذاری وابسته به
سرعت)، آفست را کم کرد، سلول‌ها را با هم‌پوشانی بیشتر در امتداد خط نصب کرد یا از تحویل شرطی استفاده کرد.
\end{solutionbox}

% ---------------------------------------------------------------------
\section{مطالعهٔ بیشتر}
تحرک در \acro{LTE} در \cite{ts36300}، رویه‌های حالت بیکار در \cite{ts36304} و رویدادهای اندازه‌گیری در
\cite{ts36331} توصیف شده‌اند؛ تحویل \lr{GSM/UMTS} در \cite{ts23009} و تحرک \acro{NR}، \acro{CHO} و
\acro{DAPS} در \cite{ts38300}. \lr{Akyildiz} و همکاران \cite{akyildiz1999mobility} مدیریت مکان و مدل جریان
سیال را مرور می‌کنند. \lr{Sesia} و همکاران \cite{sesia2011lte} و \lr{Holma} و \lr{Toskala}
\cite{holma2011lte} تنظیم پارامترهای تحویل در \acro{LTE} را بررسی می‌کنند.

\begin{latin}
\printbibliography[heading=subbibliography,title={References}]
\end{latin}
\end{document}
